%------------------------------------------------------------------------------%
\begin{question}
  Considerando a equação
  \[
    x^2+y^2-4x+2y+6=0
  \]
  \vspace{-\baselineskip}
  \begin{enumerate}
    \item identifique a curva
    \item determine seus principais valores e pontos
    \item esboce a curva
  \end{enumerate}
\end{question}

%------------------------------------------------------------------------------%
 \begin{resolution}{Completando quadrados}
   Agrupamos os termos

   \[
     x^2-4x+y^2+2y=-6
   \]

   Completamos os quadrados

   \begin{align*}
     x^2-4x+4&=(x-2)^2
     \\
     y^2+2y+1&=(y+1)^2
   \end{align*}

   Acrescentamos os termos correspondentes

   \begin{align*}
     (x-2)^2+(y+1)^2
     &= -6+4+1
     \\
     &= -1
   \end{align*}
 \end{resolution}

 %------------------------------------------------------------------------------%
 \begin{resolution}{Classificação}
   A equação obtida é

   \[
     (x-2)^2+(y+1)^2=-1
   \]

   O lado esquerdo é uma soma de quadrados e, portanto, é
   sempre maior ou igual a zero

   Não existem valores reais de $x$ e $y$ que satisfaçam a equação

   A cônica não possui pontos reais

   Portanto, trata-se de uma cônica degenerada correspondente
   a uma circunferência imaginária
 \end{resolution}

%------------------------------------------------------------------------------%
\begin{resolution}{Esboço}
  \centering
  \includegraphics{I-07-Identificacao-ex-D}
\end{resolution}

%------------------------------------------------------------------------------%
